\chapter{Harish-Chandra 模}
\label{ch:04.4}

\section{旗簇上的群作用}
\label{sec:04.4.1}

设 $\g$ 是半单 Lie 代数，$G=\operatorname{Int}(\g)$。下面的结果在本书后面起重要作用。

\begin{proposition}\label{prop:04.4.1.1}
设 $K$ 是 $G$ 的子群。则旗簇 $X$ 中的 $K$-轨道都是仿射嵌入的。
\end{proposition}

证明基于下面几个观察。

\begin{lemma}\label{lem:04.4.1.2}
设 $S$ 是可解代数群，$S'$ 是它的闭子群。则 $S/S'$ 是仿射簇。
\end{lemma}

\begin{proof}
先设 $S$ 幂幺。设 $\s$ 是 $S$ 的 Lie 代数，$\s'$ 是 $S'$ 的 Lie 代数。若 $S'\neq S$，则存在 $\s$ 的余维 $1$ 的 Lie 子代数 $\mathfrak{r}$，它包含 $\s'$。由于指数映射是 $\s$ 到 $S$ 的同构，簇 $S$ 同构于仿射直线与由 $\mathfrak{r}$ 决定的闭子群 $R$ 的乘积。又 $S/S'$ 同构于仿射直线与 $R/S'$ 的乘积。对 $S'$ 在 $S$ 中的余维作归纳，即得 $S/S'$ 是仿射空间。

现设 $S$ 任意，$S'$ 幂幺。则 $S'$ 是 $S$ 的幂幺根 $N$ 的闭子群 [2, III.10.6]。由 Levi 分解，此时 $S/S'$ 同构于 $S$ 的极大环面 $T$ 与 $N/S'$ 的乘积。这把证明归结为第一种情形。

最后设 $S'$ 任意。设 $N'$ 是它的幂幺根，$T'$ 是 $S'$ 中的极大环面。由 Levi 分解，$S'$ 是 $N'$ 与 $T'$ 的半直积。由已证的第一部分，$S/N'$ 是仿射簇。群 $T'$ 作用在簇 $S/N'$ 上，其商为 $S/S'$。由于 $T'$ 约化，这个商是仿射簇。
\begin{annote}
底本此处把 $T'$ 印作 $T$。本段中作用在 $S/N'$ 上、商为 $S/S'$ 的群是 $S'$ 的极大环面 $T'$；$T$ 是 $S$ 的极大环面，并不作用在这个商上。故按上下文改正。[来源：底本 p.85，上下文自证]
\end{annote}
\end{proof}

现在设 $Y$ 是 $G$ 的齐性空间。定义 $G$ 在 $Y\times X$ 上的作用为
$$g(y,x)=(gy,gx),\qquad g\in G,\ x\in X,\ y\in Y.$$

\begin{lemma}\label{lem:04.4.1.3}
$Y\times X$ 中的 $G$-轨道都是仿射嵌入的。
\end{lemma}

\begin{proof}
固定一点 $v\in X$。设 $B$ 是对应于 $v$ 的 Borel 子群。$Y\times X$ 中每条 $G$-轨道都与 $Y\times\{v\}$ 相交。设 $u\in Y$。则 $(u,v)$ 的 $G$-轨道 $Q$ 与 $Y\times\{v\}$ 的交等于 $Bu\times\{v\}$。设 $\bar N$ 是与 $B$ 相对的 Borel 子群的幂幺根。则 $\bar Nv$ 是 $v$ 在 $X$ 中的开邻域，且 $\bar n\mapsto\bar nv$ 是 $\bar N$ 到这个邻域的同构。$Q$ 与 $Y\times\bar Nv$ 的交等于簇 $Bu\times\bar N$ 在映射 $(y,\bar n)\mapsto\bar n(y,v)$ 下的像，而这个映射显然是浸入。由于 $B$ 可解，由 \ref{lem:04.4.1.2}，它的轨道 $Bu$ 是仿射簇。于是 $Bu\times\bar N$ 是仿射簇，从而 $Q$ 与 $Y\times\bar Nv$ 的交是仿射的。因此可以构造 $Y\times X$ 的一个开覆盖，使 $Q$ 与其中每个成员的交都是仿射簇。由于态射的仿射性关于靶簇是局部性质，证明到此完成。
\end{proof}

现在证明 \ref{prop:04.4.1.1}。令 $Y=G/K$，$u\in Y$ 是单位陪集\index{peiji@陪集|hyperpage}。则浸入 $i_u:X\to Y\times X$（$i_u(x)=(u,x)$）的像是 $Y\times X$ 的闭子簇，且同构于 $X$。设 $v\in X$，$Q'$ 是 $v$ 的 $K$-轨道。则 $i_u$ 的像与 $Y\times X$ 中 $(u,v)$ 的 $G$-轨道 $Q$ 的交等于
$$i_u(X)\cap Q=(\{u\}\times X)\cap Q=\{u\}\times Q'.$$
设 $U$ 是 $Y\times X$ 中的开仿射子集。则 $U\cap Q$ 是 $Q$ 的开仿射子集。
\begin{annote}
底本此处引作「由 2」，即引理 1.2。但 1.2 只断言可解群的商簇仿射；此处要用的是刚证的 \ref{lem:04.4.1.3}，即 $Y\times X$ 中轨道的仿射嵌入。故按 1.3 引用。[来源：底本 p.85，上下文自证]
\end{annote}
又因 $i_u(X)$ 在 $Y\times X$ 中闭，$U\cap(i_u(X)\cap Q)$ 是 $i_u(X)\cap Q$ 的开仿射子集。这推出 $U\cap(\{u\}\times Q')$ 是 $\{u\}\times Q'$ 的开仿射子集。此外，由于 $i_u$ 是闭浸入，$V=i_u^{-1}(U)$ 是 $X$ 的开仿射子集，而 $V\cap Q'$ 是 $Q'$ 的开仿射子集。显然这推出 $Q'$ 仿射嵌入 $X$，从而完成 \ref{prop:04.4.1.1} 的证明。

\section{Harish-Chandra 对}
\label{sec:04.4.2}

设 $K$ 是代数群，$\phi:K\to\operatorname{Int}(\g)$ 是代数群态射，且其微分是单射。此时可把 $K$ 的 Lie 代数与 $\g$ 的子代数 $\kk$ 等同。显然，$K$ 自然作用在 $X$ 上。

若 $K$ 在 $X$ 上的作用只有有限多条轨道，则称对 $(\g,K)$ 是 Harish-Chandra 对。

若 $(\g,K)$ 是 Harish-Chandra 对，则 $K$ 在 $X$ 中有开轨道。事实上，这两个性质等价 [?]。

\begin{theorem}\label{thm:04.4.2.1}
设 $K$ 是 $\operatorname{Int}(\g)$ 的闭子群。则下列条件等价：
\begin{enumerate}[label=(\roman*)]
\item $K$ 在 $X$ 中有开轨道；
\item $K$ 在 $X$ 中只有有限多条轨道。
\end{enumerate}
\end{theorem}

Harish-Chandra 对的一个例子是 $(\g,B)$，其中 $B$ 是 $\operatorname{Int}(\g)$ 的 Borel 子群。$B$ 在 $X$ 中轨道的有限性就是 Bruhat 引理 [2, 14.11]。

另一类重要的例子如下。设 $\sigma$ 是 $\g$ 的对合，$\kk$ 是 $\g$ 中被 $\sigma$ 固定的全体向量构成的 Lie 子代数。称 $\kk$ 是 $\g$ 的对合子代数。

\begin{proposition}\label{prop:04.4.2.2}
设 $K$ 是 $\operatorname{Int}(\g)$ 的闭子群，其 Lie 代数 $\kk$ 是 $\g$ 的对合子代数。则 $K$ 在 $X$ 上的作用只有有限多条轨道。
\end{proposition}

把 $G=\operatorname{Int}(\g)$ 的微分为 $\sigma$ 的对合自同构仍记作 $\sigma$。证明的关键是下面的引理。先在 $X\times X$ 上定义 $G$ 的作用：
$$g(x,y)=(gx,\sigma(g)y),\qquad g\in G,\ x,y\in X.$$

\begin{lemma}\label{lem:04.4.2.3}
群 $G$ 在 $X\times X$ 上的作用只有有限多条轨道。
\end{lemma}

\begin{proof}
固定一点 $v\in X$。设 $B_v$ 是对应于 $v$ 的 $G$ 的 Borel 子群，并令 $B=\sigma(B_v)$。$X\times X$ 中每条 $G$-轨道都与 $X\times\{v\}$ 相交。设 $u\in X$。则过 $(u,v)$ 的 $G$-轨道 $Q$ 与 $X\times\{v\}$ 的交等于 $Bu\times\{v\}$。由 Bruhat 分解 [2, IV.14.11]，这推出 $X\times X$ 中 $G$-轨道的条数有限。
\end{proof}

现在说明 \ref{prop:04.4.2.2} 是 \ref{lem:04.4.2.3} 的推论。先可设 $K$ 连通。设 $\Delta$ 是 $X\times X$ 的对角线。由 \ref{lem:04.4.2.3}，$X\times X$ 的轨道分层诱导 $\Delta$ 的一个分层，其层是有限多个不可约的仿射嵌入子簇，它们正是各 $G$-轨道与 $\Delta$ 的交的不可约分支。这些层 $K$-不变，从而是若干 $K$-轨道的并。设 $V$ 是其中一者，$(x,x)\in V$，$Q$ 是 $(x,x)$ 的 $K$-轨道。设 $\bb_x$ 是对应于 $x$ 的 $\g$ 的 Borel 子代数，则 $X$ 在 $x$ 处的切空间 $T_x(X)$ 可等同于 $\g/\bb_x$。设 $p_x$ 是 $\g$ 到 $\g/\bb_x$ 的投影。$X\times X$ 在 $(x,x)$ 处的切空间 $T_{(x,x)}(X\times X)$ 可等同于 $\g/\bb_x\times\g/\bb_x$。若轨道映射 $f:G\to X\times X$ 由 $f(g)=g(x,x)$ 定义，则它在 $G$ 的单位元处的微分是线性映射 $\xi\mapsto(p_x(\xi),p_x(\sigma(\xi)))$，从 $\g$ 映入 $\g/\bb_x\times\g/\bb_x$。于是 $V$ 在 $(x,x)$ 处的切空间含于这个微分的像与 $T_{(x,x)}(X\times X)$ 中对角线的交，即
\begin{align*}
T_{(x,x)}(V)&\subset\{(p_x(\xi),p_x(\xi))\mid \xi\in\g,\ p_x(\xi)=p_x(\sigma(\xi))\}\\&=\{(p_x(\xi),p_x(\xi))\mid \xi\in\kk\}=T_{(x,x)}(Q).
\end{align*}
于是 $V$ 在 $(x,x)$ 处的切空间与 $Q$ 的切空间相同，而 $Q$ 在 $V$ 中开。由 $V$ 不可约，这推出 $V$ 就是一条 $K$-轨道。因此对角线 $\Delta$ 的这个分层正是 $X$ 的 $K$-轨道分层经对角映射诱导的分层。于是 \ref{prop:04.4.2.2} 得证。

下面结果只是 \ref{prop:04.4.2.2} 的改写。

\begin{theorem}\label{thm:04.4.2.4}
设 $K$ 是代数群，$\phi:K\to\operatorname{Int}(\g)$ 是代数群态射，且其微分单射。设 $\g$ 的 Lie 子代数 $\kk$ 是对合子代数。则 $(\g,K)$ 是 Harish-Chandra 对。
\end{theorem}

这样的 Harish-Chandra 对称为对合 Harish-Chandra 对。

\section{Harish-Chandra 模与 Harish-Chandra 层}
\label{sec:04.4.3}

设 $(\g,K)$ 是 Harish-Chandra 对。Harish-Chandra 模 $V$ 是一个向量空间，满足：
\begin{enumerate}[label=(\roman*)]
\item $V$ 是有限生成 $U(\g)$-模，且作为 $Z(\g)$-模是局部有限的；
\item $V$ 是代数 $K$-模；
\item $\g$ 与 $K$ 的作用相容，即由 $K$-作用的微分给出的 $\kk$-作用与 $\kk$ 作为 $\g$ 的子代数所给的作用相同，并且
$$(\phi(k)\xi)\cdot v=k\cdot\xi\cdot k^{-1}\cdot v$$
对 $k\in K$、$\xi\in\g$、$v\in V$ 成立。
\end{enumerate}

Harish-Chandra 模的态射是既为 $U(\g)$-模态射又为 $K$-模态射的线性映射。把 Harish-Chandra 模构成的范畴记作 $\mathcal{M}_{fg}(U(\g),K)$。对 $\lambda\in\h^*$、$\theta=W\cdot\lambda$，把 $\mathcal{M}_{fg}(U(\g),K)$ 中无穷小特征标为 $\chi_\lambda$ 的模构成的全子范畴记作 $\mathcal{M}_{fg}(U_\theta,K)$。

$\mathcal{M}_{coh}(\D_\lambda,K)$ 的对象称为 Harish-Chandra 层。

设 $\lambda\in\h^*$，$\theta=W\cdot\lambda$。由 ……，显然对 $\mathcal{M}_{fg}(U_\theta,K)$ 的任一对象 $V$，其局部化 $\Delta_\lambda(V)$ 是 $\mathcal{M}_{coh}(\D_\lambda,K)$ 的对象。又由 L.1.21，$\mathcal{M}_{coh}(\D_\lambda,K)$ 中 Harish-Chandra 层 $\mathcal{V}$ 的上同调模 $H^i(X,\mathcal{V})$（$0\le i\le\dim X$）作为 $U_\theta$-模是有限生成的。由 ……，它们又是代数 $K$-模，故属于 $\mathcal{M}_{fg}(U_\theta,K)$。

\begin{lemma}\label{lem:04.4.3.1}
任一 Harish-Chandra 层 $\mathcal{V}$ 都有良滤过 $F\mathcal{V}$，其中每一项都是 $K$-齐性凝聚 $\calO_X$-模。
\end{lemma}

\begin{proof}
用 $\calO(\mu)$（$\mu\in P(\Sigma)$ 充分负）作扭，可设 $\lambda$ 反支配且正则。此时由 L.1.3，$\mathcal{V}=\D_\lambda\otimes_{U_\theta}V$，其中 $V=\Gamma(X,\mathcal{V})$。由于 $V$ 是代数 $K$-模，也是有限生成 $U_\theta$-模，存在有限维 $K$-不变子空间 $U$，它在 $U_\theta$ 下生成 $V$。于是 $F_p\D_\lambda\otimes_\C U$（$p\in\Z_+$）都是 $K$-齐性凝聚 $\calO_X$-模。自然映射 $F_p\D_\lambda\otimes_\C U\to\mathcal{V}$ 是 $K$-等变的，故其像 $F_p\mathcal{V}$ 对任意 $p\in\Z_+$ 都是 $\mathcal{V}$ 的 $K$-齐性凝聚 $\calO_X$-子模。

我们断言 $F\mathcal{V}$ 是 $\D_\lambda$-模 $\mathcal{V}$ 的良滤过。显然它是 $\mathcal{V}$ 的 $\D_\lambda$-模滤过，且各项都是 $K$-齐性凝聚 $\calO_X$-模。由于 $\mathcal{V}$ 由其整体截面生成，要证滤过穷尽，只需证 $\mathcal{V}$ 的任一整体截面 $v$ 对充分大的 $p$ 落在 $F_p\mathcal{V}$ 中。由于 $V$ 作为 $U_\theta$-模由 $U$ 生成，存在 $T_i\in U_\theta$、$u_i\in U$（$1\le i\le m$）使 $v=\sum_{i=1}^m T_iu_i$。另一方面，存在 $p\in\Z_+$，使 $T_i$（$1\le i\le m$）都是 $F_p\D_\lambda$ 的整体截面。这推出 $v\in F_p\mathcal{V}$。最后，由 $F\mathcal{V}$ 的构造显然有 $F_p\D_\lambda F_q\mathcal{V}=F_{p+q}\mathcal{V}$ 对一切 $p,q\in\Z_+$ 成立，即 $F\mathcal{V}$ 是良滤过。
\end{proof}

关于 Harish-Chandra 层，关键结果如下。

\begin{theorem}\label{thm:04.4.3.2}
Harish-Chandra 层是完整 $\D_\lambda$-模。特别地，它们长度有限。
\end{theorem}

我们实际上将证明更强的结果。先作一些记号。设 $Y$ 是纯维数的光滑代数簇，$Z$ 是 $Y$ 的光滑子簇。在 $T^*(Y)$ 中定义光滑子簇 $N_Z(Y)$，它是所有满足 $z\in Z$、且 $\omega\in T_z(Y)^*$ 在 $T_z(Z)\subset T_z(Y)$ 上为零的点 $(z,\omega)\in T^*(Y)$ 构成的簇。称 $N_Z(Y)$ 为 $Z$ 在 $Y$ 中的余法簇（conormal variety）。

\begin{lemma}\label{lem:04.4.3.3}
$Z$ 在 $Y$ 中的余法簇 $N_Z(Y)$ 的维数等于 $\dim Y$。
\end{lemma}

\begin{proof}
在 $T_z(Y)^*$ 中，在 $T_z(Z)$ 上为零的全体线性形式构成的空间维数等于 $\dim T_z(Y)-\dim T_z(Z)=\dim Y-\dim_z Z$。故 $\dim_z N_Z(Y)=\dim Y$。
\end{proof}

设 $\lambda\in\h^*$。由 ……，$\Gr\D_\lambda=\pi_*(\calO_{T^*(X)})$，其中 $\pi:T^*(X)\to X$ 是自然投影。设 $\xi\in\g$。则 $\xi$ 决定 $\D_\lambda$ 的一个阶 $\le 1$ 的整体截面，即 $F_1\D_\lambda$ 的整体截面。因此该截面的象征\index{xiangzheng@象征|hyperpage}是 $\pi_*(\calO_{T^*(X)})$ 的整体截面，且与 $\lambda$ 无关。设 $x\in X$。则轨道映射 $f_x:G\to X$（$f_x(g)=gx$）在 $1\in G$ 处的微分把 Lie 代数 $\g$ 满射到 $x$ 处的切空间 $T_x(X)$。这个映射的核是 $\bb_x$，即 $f_x$ 在 $1$ 处的微分 $T_1(f_x)$ 把 $\g/\bb_x$ 与 $T_x(X)$ 等同。由 $\xi$ 决定的截面的象征由函数
$$(x,\omega)\mapsto\omega(T_1(f_x)(\xi)),\qquad x\in X,\ \omega\in T_x(X)^*,$$
给出。

设 $K$ 是 $\operatorname{Int}(\g)$ 的闭子群，$\kk$ 是它的 Lie 代数。把 $\calO_X$-模 $\pi_*(\calO_{T^*(X)})$ 中由 $\kk$ 的元素所对应截面的象征生成的理想记作 $I_K$。设 $N_K$ 是这个理想在 $T^*(X)$ 中的零点集。

\begin{lemma}\label{lem:04.4.3.4}
簇 $N_K$ 是 $X$ 中一切 $K$-轨道 $Q$ 的余法簇 $N_Q(X)$ 的并。
\end{lemma}

\begin{proof}
设 $x\in X$，$Q$ 是过 $x$ 的 $K$-轨道。则
\begin{align*}
N_K\cap T_x(X)^*&=\{\omega\in T_x(X)^*\mid \omega\ \text{在}\ T_1(f_x)(\kk)\ \text{上为零}\}\\
&=\{\omega\in T_x(X)^*\mid \omega\ \text{在}\ T_x(Q)\ \text{上为零}\}=N_Q(X)\cap T_x(X)^*,
\end{align*}
即 $N_K$ 是全体 $N_Q(X)$ 的并。
\end{proof}

\begin{corollary}\label{cor:04.4.3.5}
设 $K$ 在 $X$ 上的作用只有有限多条轨道。则：
\begin{enumerate}[label=(\roman*)]
\item $\dim N_K=\dim X$；
\item 若 $K$ 连通，则 $N_K$ 的不可约分支是 $X$ 中各 $K$-轨道 $Q$ 的余法簇 $N_Q(X)$ 的闭包 $\overline{N_Q(X)}$。
\end{enumerate}
\end{corollary}

\begin{proof}
对 $X$ 中任一 $K$-轨道 $Q$，由 \ref{lem:04.4.3.3}，其余法簇 $N_Q(X)$ 的维数等于 $\dim X$。由于 $X$ 中 $K$-轨道只有有限多条，由 \ref{lem:04.4.3.3} 与 \ref{lem:04.4.3.4}，$N_K$ 是有限多个维数为 $\dim X$ 的子簇的并。这推出 (i)。又
$$N_K=\bigcup_Q N_Q(X).$$
若 $K$ 连通，它在 $X$ 中的轨道也连通，故余法簇 $N_Q(X)$ 也连通。这些簇光滑，于是立即推出它们不可约。因此它们的闭包 $\overline{N_Q(X)}$ 是 $N_K$ 中维数为 $\dim X$ 的不可约闭子簇，从而是 $N_K$ 的不可约分支。这证明了 (ii)。
\end{proof}

因此，\ref{thm:04.4.3.2} 是下面结果的直接推论。

\begin{proposition}\label{prop:04.4.3.6}
设 $\mathcal{V}$ 是 Harish-Chandra 层。则 $\mathcal{V}$ 的特征簇\index{tezhengcu@特征簇|hyperpage} $\operatorname{Char}(\mathcal{V})$ 是 $N_K$ 的闭子簇。
\end{proposition}

\begin{proof}
由 \ref{lem:04.4.3.1}，$\mathcal{V}$ 有良滤过 $F\mathcal{V}$，其各项是 $K$-齐性凝聚 $\calO_X$-模。因此 $\D_\lambda$ 中对应于 $\kk$ 的整体截面把 $F_p\mathcal{V}$ 映入自身（$p\in\Z$）。于是它们的象征零化 $\Gr\mathcal{V}$，从而 $I_K$ 含于 $\Gr\mathcal{V}$ 在 $\pi_*(\calO_{T^*(X)})$ 中的零化子。这推出特征簇 $\operatorname{Char}(\mathcal{V})$ 是 $N_K$ 的闭子簇。
\end{proof}

下面结果是 \ref{thm:04.4.3.2} 的直接推论。

\begin{theorem}\label{thm:04.4.3.7}
每个 Harish-Chandra 模都是有限长的。
\end{theorem}

\begin{proof}
设 $V$ 是 Harish-Chandra 模。由于它作为 $U(\g)$-模有限生成、作为 $Z(\g)$-模局部有限，存在有限维 $Z(\g)$-子模 $U$，它在 $U(\g)$ 下生成 $V$。因此存在有限集 $\lambda_1,\lambda_2,\dots,\lambda_k\in\h^*$ 与 $n\in\N$，使
$$P(\xi)=\prod_{i=1}^k(\xi-\chi_{\lambda_i}(\xi))^n,\qquad \xi\in Z(\g),$$
零化 $U$。由于 $U$ 作为 $U(\g)$-模生成 $V$，推出对每个 $\xi\in Z(\g)$，$P(\xi)$ 零化 $V$。因此 $V$ 是子模 $V_i$（$1\le i\le k$）的直和，其中
$$V_i=\{v\in V\mid (\xi-\chi_{\lambda_i}(\xi))^nv=0\ \text{对一切}\ \xi\in Z(\g)\ \text{与充分大的}\ n\in\N\}.$$

因此可设 $V$ 被 $(\xi-\chi_\lambda(\xi))^n$（$\xi\in Z(\g)$）零化。我们断言这样的 $V$ 有有限滤过 $FV$，其各项是 Harish-Chandra 子模，且相应的分次模具有无穷小特征标 $\chi_\lambda$。对 $\dim U$ 作归纳。若 $\dim U=1$，断言显然。一般地，$U$ 含 $Z(\g)$ 的一维特征子空间 $U_1$。显然 $U_1$ 生成 $V$ 的一个 Harish-Chandra 子模 $V_1$，其无穷小特征标为 $\chi_\lambda$。商 $V_2=V/V_1$ 由 $U$ 的像 $U_2$ 生成，且 $\dim U_2\le\dim U-1$。

因此可设 $V$ 具有无穷小特征标，即 $V\in\mathcal{M}_{fg}(U_\theta,K)$。此时可在 $\theta$ 中取反支配的 $\lambda$。局部化 $\Delta_\lambda(V)$ 是 Harish-Chandra 层，故由 \ref{thm:04.4.3.2} 长度有限。由 L.1.1，$V=\Gamma(X,\Delta_\lambda(V))$。因此 $\Gamma$ 的正合性与 L.4.1 推出 $V$ 长度有限。
\end{proof}

\section{Harish-Chandra 模的 $n$-同调}
\label{sec:04.4.4}

设 $V$ 是 Harish-Chandra 模。对任意 $x\in X$，$V$ 的 $\n_x$-同调可由标准复形 $C_\bullet(\n_x,V)$ 计算：
$$C_p(\n_x,V)=\wedge^{-p}\n_x\otimes_\C V,\qquad p\in\Z,$$
其微分为
\begin{align*}
d(\xi_1\wedge\xi_2\wedge\cdots\wedge\xi_p\otimes v)
={}&\sum_{i=1}^{p}(-1)^{i+1}\xi_1\wedge\cdots\wedge\hat\xi_i\wedge\cdots\wedge\xi_p\otimes\xi_i v\\
&+\sum_{i<j}(-1)^{i+j}[\xi_i,\xi_j]\wedge\xi_1\wedge\cdots\wedge\hat\xi_i\wedge\cdots\wedge\hat\xi_j\wedge\cdots\wedge\xi_p\otimes v,
\end{align*}
其中 $\xi_1,\dots,\xi_p\in\n_x$，$v\in V$。

这立即推出，对 $X$ 中同一 $K$-轨道 $Q$ 上的各点 $x$，$\h$-模 $H_p(\n_x,V)$ 典范同构。

\begin{theorem}\label{thm:04.4.4.1}
设 $V$ 是 Harish-Chandra 模，$x\in X$。则一切 $\n_x$-同调模 $H_p(\n_x,V)$（$p\in\Z$）都是有限维的。
\end{theorem}

\begin{proof}
在 \ref{thm:04.4.3.7} 的证明中，我们构造了 $V$ 的一个有限滤过，其各项是 Harish-Chandra 子模，其合成因子是具有无穷小特征标的 Harish-Chandra 模。因此，用滤过对象的谱序列（……），立见只需对 $V\in\mathcal{M}_{fg}(U_\theta,K)$ 证明结论。

若 $\theta$ 是 $W$ 的正则轨道，结论由 L.5.4 与 \ref{thm:04.4.3.2} 得出。现设 $\theta$ 是任意的 Weyl 群轨道。固定反支配的 $\lambda\in\theta$。设 $F$ 是最高权 $\mu\in P(\Sigma)$ 为正则的有限维不可约 $\g$-模。则 $\lambda-\mu$ 正则反支配，且由 C.2.1 有
$$\Delta_\lambda(V)=(\Delta_\lambda(V)(-\mu)\otimes_{\calO_X}F)[\lambda].$$
因此
$$V=\Gamma(X,\Delta_\lambda(V))=(\Gamma(X,\Delta_\lambda(V)(-\mu))\otimes_\C F)[\lambda].$$

显然 $U=\Gamma(X,\Delta_\lambda(V)(-\mu))$ 是无穷小特征标 $\chi_{\lambda-\mu}$ 正则的 Harish-Chandra 模，且 $H_p(\n_x,U)$（$p\in\Z$）有限维。设 $(F_i;1\le i\le m)$ 是 $F$ 作为 $\n_x$-模的递增 Jordan–Hölder 滤过。则对应的分次模是平凡 $\n_x$-模。因此 $(U\otimes_\C F_i;1\le i\le m)$ 是 $U\otimes_\C F$ 的 $\n_x$-模滤过，其对应的分次模是 $m$ 个 $U$ 的直和。由滤过对象（……）的谱序列推出 $H_p(\n_x,U\otimes F)$（$p\in\Z$）有限维。由 $V=(U\otimes_\C F)[\lambda]$ 即得结论。
\end{proof}

下面结果深刻得多，它来自 L.8 的主要结果。设 $Q^o$ 是 $X$ 中唯一的 $K$-开轨道。

\begin{theorem}\label{thm:04.4.4.2}
设 $V$ 是 Harish-Chandra 模。若 $V$ 非零，则对一切 $x\in Q^o$ 有 $H_0(\n_x,V)\neq 0$。
\end{theorem}

\begin{proof}
由于 $Q^o$ 在 $X$ 中开且稠密，且 $H_0(\n_x,V)$ 对一切 $x\in Q^o$ 典范同构，看出 $\lambda\in\h^*$ 是 Harish-Chandra 模 $V$ 的 exponent 当且仅当 $H_0(\n_x,V)(\lambda+\rho)\neq 0$（$x\in Q^o$）。于是结论由 L.8.11 得出。
\begin{annote}
底本此处印作 $\lambda\in\h$。但 $H_0(\n_x,V)(\lambda+\rho)$ 是 $\h^*$ 中的权，且本章前文一律写 $\lambda\in\h^*$，故应为 $\lambda\in\h^*$。已改。[来源：底本 p.91，上下文自证]
\end{annote}
\end{proof}

\section{不可约 Harish-Chandra 层}
\label{sec:04.4.5}

现在描述 Harish-Chandra 对 $(\g,K)$ 的全部不可约 Harish-Chandra 层。为简单起见，设 $K$ 连通。先看下面的注记。

\begin{lemma}\label{lem:04.4.5.1}
设 $\mathcal{V}$ 是不可约 Harish-Chandra 层。则它的支集 $\operatorname{supp}(\mathcal{V})$ 是 $X$ 中某条 $K$-轨道 $Q$ 的闭包。
\end{lemma}

\begin{proof}
由于 $K$ 连通，Harish-Chandra 层 $\mathcal{V}$ 不可约当且仅当它作为 $\D_\lambda$-模不可约。为看出这一点，可用 $\calO(\mu)$（$\mu$ 充分负）作扭，从而设 $\lambda$ 反支配且正则。此时结论由范畴等价与 Harish-Chandra 模的相应结论（它显然成立）得出。

因此由 ……，$\operatorname{supp}(\mathcal{V})$ 是 $X$ 的不可约闭子簇。它又必须 $K$-不变，故是若干 $K$-轨道的并。$K$-轨道只有有限多条，故 $\operatorname{supp}(\mathcal{V})$ 中存在轨道 $Q$，使 $\dim Q=\dim\operatorname{supp}(\mathcal{V})$。于是 $\overline Q$ 是 $\operatorname{supp}(\mathcal{V})$ 的不可约闭子集，且 $\dim\overline Q=\dim\operatorname{supp}(\mathcal{V})$。这推出 $\overline Q=\operatorname{supp}(\mathcal{V})$。
\end{proof}

设 $\mathcal{V}$ 是不可约 Harish-Chandra 层，$Q$ 是 $X$ 中使 $\operatorname{supp}(\mathcal{V})=\overline Q$ 的 $K$-轨道。令 $X'=X-\partial Q$。则 $X'$ 是 $X$ 的开子簇，$Q$ 是 $X'$ 的闭子簇。由 ……，$\mathcal{V}$ 在 $X'$ 上的限制 $\mathcal{V}|_{X'}$ 不可约。设 $i:Q\to X$、$i':Q\to X'$、$j:X'\to X$ 是自然浸入。于是 $i=j\circ i'$。因此 $Ri^!=R(i')^!\circ j^!$，其中 $j^!=j^+$ 就是到 $X$ 的开子簇 $X'$ 的通常限制。由此对 $p\in\Z$ 有 $R^pi^!(\mathcal{V})=R^p(i')^!(\mathcal{V}|_{X'})$。

由于 $i':Q\to X'$ 是闭光滑子簇的浸入，且 $\operatorname{supp}(\mathcal{V}|_{X'})=Q$，由 Kashiwara 范畴等价，$R^pi^!(\mathcal{V})=0$（$p\neq 0$），且 $\tau=i^!(\mathcal{V})$ 是 $Q$ 上的不可约 $(\D^i_\lambda,K)$-模。此外 $(i')_+(\tau)=\mathcal{V}|_{X'}$。由 \ref{thm:04.4.3.2}，$\mathcal{V}$ 完整，故 $\tau$ 是完整模。由 ……，存在 $Q$ 的开稠密子集 $U$，使 $\tau|_U$ 是联络。由于 $K$ 在 $Q$ 上可迁，$\tau$ 必是 $K$-齐性联络。

因此，对每个不可约 Harish-Chandra 层，都附有一个对 $(Q,\tau)$，它由 $K$-轨道 $Q$ 与 $Q$ 上不可约的 $K$-齐性联络 $\tau$ 组成，并满足：
\begin{enumerate}[label=(\roman*)]
\item $\operatorname{supp}(\mathcal{V})=\overline Q$；
\item $i^!(\mathcal{V})=\tau$。
\end{enumerate}
称对 $(Q,\tau)$ 为附于 $\mathcal{V}$ 的标准数据。

设 $Q$ 是 $X$ 中的 $K$-轨道，$\tau\in\mathcal{M}(\D^i_\lambda,K)$ 是 $Q$ 上不可约的 $K$-齐性联络。则由 ……，$I(Q,\tau)=i_+(\tau)$ 是 $(\D_\lambda,K)$-模。又由 ……，它完整，即 $I(Q,\tau)$ 是 Harish-Chandra 层。称它为附于 $(Q,\tau)$ 的标准 Harish-Chandra 层。

\begin{lemma}\label{lem:04.4.5.2}
设 $Q$ 是 $X$ 中的 $K$-轨道，$\tau$ 是 $Q$ 上不可约的 $K$-齐性联络。则标准 Harish-Chandra 层 $I(Q,\tau)$ 含唯一的不可约 Harish-Chandra 子层。
\end{lemma}

\begin{proof}
显然
$$I(Q,\tau)=i_+(\tau)=j_+((i')_+(\tau)).$$
因此 $I(Q,\tau)$ 没有支集含于 $\partial Q$ 的截面。于是 $I(Q,\tau)$ 的任一非零 $\D_\lambda$-子模 $U$ 在 $X'$ 上的限制非零。由 Kashiwara 范畴等价，$(i')_+(\tau)$ 是不可约 $(\D_\lambda|_{X'})$-模。故 $U|_{X'}=I(Q,\tau)|_{X'}$。因此对 $I(Q,\tau)$ 的任意两个非零 $\D_\lambda$-子模 $U$ 与 $U'$，有 $U\cap U'\neq 0$。由于 $I(Q,\tau)$ 长度有限，它有极小 $\D_\lambda$-子模；由前一段，这个子模唯一。由唯一性，它必 $K$-等变，从而是 Harish-Chandra 层。
\end{proof}

把 $I(Q,\tau)$ 的唯一的不可约 Harish-Chandra 子层记作 $L(Q,\tau)$。下面的结果给出不可约 Harish-Chandra 层的分类。

\begin{theorem}\label{thm:04.4.5.3}
\begin{enumerate}[label=(\roman*)]
\item 具标准数据 $(Q,\tau)$ 的不可约 Harish-Chandra 层 $\mathcal{V}$ 同构于 $L(Q,\tau)$。
\item 设 $Q$、$Q'$ 是 $X$ 中的 $K$-轨道，$\tau$、$\tau'$ 分别是 $Q$、$Q'$ 上的不可约 $K$-齐性联络。则 $L(Q,\tau)\cong L(Q',\tau')$ 当且仅当 $Q=Q'$ 且 $\tau\cong\tau'$。
\end{enumerate}
\end{theorem}

\begin{proof}
(i) 设 $\mathcal{V}$ 是不可约 Harish-Chandra 层，$(Q,\tau)$ 是相应的标准数据。则如前所述，$\mathcal{V}|_{X'}=(i')_+(\tau)$。由 $j_+$ 的泛性质，存在 $\mathcal{V}$ 到 $I(Q,\tau)=j_+((i')_+(\tau))$ 的非平凡态射，它延拓这个同构。由于 $\mathcal{V}$ 不可约，其核必为零；由 \ref{lem:04.4.5.2}，其像必为 $L(Q,\tau)$。

(ii) 由于 $Q=\operatorname{supp}L(Q,\tau)$，显然 $L(Q,\tau)\cong L(Q',\tau')$ 推出 $Q=Q'$。其余由公式 $\tau=i^!(L(Q,\tau))$ 得出。
\end{proof}
